\section*{Sets and Indices}

Let the planning years be indexed by
\[
t \in \{1,2,3\}.
\]
Because the implementation uses distinct variables for each year-specific investment channel, the formulation is written directly by year rather than with a more abstract indexed family.

\section*{Parameters}

All parameter values are read from \texttt{general\_parameters.csv}.

\[
I \quad \text{(code: \texttt{initial\_fund})}
\]
Initial fund available at the start of year 1.

\[
r \quad \text{(code: \texttt{annual\_rate})}
\]
Annual return rate for the standard one-year investment.

\[
L_1 \quad \text{(code: \texttt{inv\_limit\_y1})}
\]
Maximum amount allowed in the special investment started in year 1.

\[
R_{12} \quad \text{(code: \texttt{return\_rate\_y1y2})}
\]
Gross return multiplier for the special investment started in year 1 and recovered at the end of year 2.

\[
L_2 \quad \text{(code: \texttt{inv\_limit\_y2})}
\]
Maximum amount allowed in the special investment started in year 2.

\[
R_{23} \quad \text{(code: \texttt{return\_rate\_y2y3})}
\]
Gross return multiplier for the special investment started in year 2 and recovered at the end of year 3.

\[
L_3 \quad \text{(code: \texttt{inv\_limit\_y3})}
\]
Maximum amount allowed in the special investment started in year 3.

\[
r_3 \quad \text{(code: \texttt{profit\_rate\_y3})}
\]
Profit rate for the special investment started in year 3 and recovered one year later, so its gross multiplier is \(1+r_3\).

\[
F_1 \quad \text{(code: \texttt{fee\_y1})}, \qquad
F_2 \quad \text{(code: \texttt{fee\_y2})}, \qquad
F_3 \quad \text{(code: \texttt{fee\_y3})}
\]
Fixed management fees charged if the year-1, year-2, or year-3 special investment channel is used.

\section*{Decision Variables}

\[
a_1 \ge 0 \quad \text{(code: \texttt{a1})}
\]
Amount placed at the start of year 1 in the standard one-year investment.

\[
a_2 \ge 0 \quad \text{(code: \texttt{a2})}
\]
Amount placed at the start of year 2 in the standard one-year investment.

\[
a_3 \ge 0 \quad \text{(code: \texttt{a3})}
\]
Amount placed at the start of year 3 in the standard one-year investment.

\[
b_1 \ge 0 \quad \text{(code: \texttt{b1})}
\]
Amount placed in the year-1 special investment, with upper bound \(L_1\).

\[
c_2 \ge 0 \quad \text{(code: \texttt{c2})}
\]
Amount placed in the year-2 special investment, with upper bound \(L_2\).

\[
d_3 \ge 0 \quad \text{(code: \texttt{d3})}
\]
Amount placed in the year-3 special investment, with upper bound \(L_3\).

\[
y_{b1} \in \{0,1\} \quad \text{(code: \texttt{y\_b1})}
\]
Binary indicator equal to 1 if the year-1 special investment channel is used.

\[
y_{c2} \in \{0,1\} \quad \text{(code: \texttt{y\_c2})}
\]
Binary indicator equal to 1 if the year-2 special investment channel is used.

\[
y_{d3} \in \{0,1\} \quad \text{(code: \texttt{y\_d3})}
\]
Binary indicator equal to 1 if the year-3 special investment channel is used.

\section*{Objective}

The code maximizes final wealth at the end of year 3 net of fixed management fees:
\[
\max \; Z
\]
with
\[
Z=
\underbrace{\left(I-a_1-b_1\right)+a_1(1+r)-a_2-c_2+a_2(1+r)+b_1R_{12}-a_3-d_3}_{\text{cash\_year3}}
+ a_3(1+r) + c_2R_{23} + d_3(1+r_3)
- F_1y_{b1} - F_2y_{c2} - F_3y_{d3}.
\]

\section*{Constraints}

Special-investment activation constraints:
\[
b_1 \le L_1 y_{b1},
\]
\[
c_2 \le L_2 y_{c2},
\]
\[
d_3 \le L_3 y_{d3}.
\]

Year-1 budget:
\[
a_1 + b_1 \le I.
\]

Year-2 budget:
\[
a_2 + c_2 \le (I-a_1-b_1)+a_1(1+r).
\]

Year-3 budget:
\[
a_3 + d_3 \le (I-a_1-b_1)+a_1(1+r)-a_2-c_2+a_2(1+r)+b_1R_{12}.
\]

Variable domains:
\[
a_1,a_2,a_3,b_1,c_2,d_3 \ge 0,
\]
\[
y_{b1},y_{c2},y_{d3}\in\{0,1\}.
\]

\section*{Notes on Implementation}

The formulation above matches the implemented mixed-integer linear model in \texttt{current\_heuristic.py}.

The standard investments \(a_1,a_2,a_3\) are modeled as one-year placements earning multiplier \(1+r\). The special investments \(b_1,c_2,d_3\) are year-specific channels with limits \(L_1,L_2,L_3\) and gross return multipliers \(R_{12},R_{23},1+r_3\), respectively.

The fixed management fees are modeled exactly as in the code: each fee is subtracted in the objective if and only if the corresponding special-investment variable is positive, enforced through the binary variables and the big-\(M\) constraints \(b_1 \le L_1 y_{b1}\), \(c_2 \le L_2 y_{c2}\), and \(d_3 \le L_3 y_{d3}\).

No fee is charged to the year-by-year budget constraints in the implementation; the fees appear only as deductions in the final objective. The year-3 objective expression also follows the code literally by carrying forward intermediate cash balances and adding the matured year-3 investments separately. The model is solved with a MIP solver (\texttt{GUROBI\_CMD} through PuLP compatibility).
